%% kapa: exam=1 round=30 year=2019 period=1 form=A subject=civil_law subject_ko=민법 seq=6 number=6
%% source: Q-Net 감정평가사 기출문제 (cst00309) artlSeq=5205205 file=extracted/1008/5205205/2019년 제30회 감정평가사 1차시험 문제지/2019년 제30회 감정평가사 1차시험 문제지/2019년도 제30회 감정평가사 1차 1교시 A형.hwp
%% conversion: kordoc 4.24.51 -> tools/gen1.py (auto); review=proofread-sonnet; flags=none
\begin{question}{6}
어부 甲은 2015년 7월 1일 조업 중 태풍으로 인하여 선박이 침몰하여 실종된 후 2017년 10월 1일 실종선고를 받았다. 이 사안에 관한 설명으로 옳은 것은? (다툼이 있으면 판례에 따름)
\begin{choices}
  \item 위 실종선고를 위해 필요한 실종기간은 1년이다.
  \item 甲은 2017년 10월 1일에 사망한 것으로 간주된다.
  \item 1순위 상속인이 있더라도 2순위 상속인은 위 실종선고를 신청할 수 있다.
  \item 甲이 극적으로 살아서 종래의 주소지로 돌아오면 위 실종선고는 자동으로 취소된다.
  \item 甲의 생환으로 실종선고가 취소되면 甲의 상속인은 악의인 경우에만 상속재산을 甲에게 반환할 의무가 있다.
\end{choices}
\end{question}
\answer{1}
